A análise de ferimentos por arma de fogo (Entrada e Saída) e a estimativa da distância do disparo (MLSD - Medical-Legal Shooting Distance) são tarefas periciais desafiadoras, geralmente demandando especialistas e dependendo de condições sub-ideais de captura de imagens em cenas de crime. O advento do Aprendizado Profundo (Deep Learning) trouxe novas oportunidades para automatizar e auxiliar essas tarefas, mas ainda enfrenta dificuldades com bases de dados limitadas e desbalanceadas.
Este trabalho investiga o uso de uma arquitetura híbrida denominada TransferKAN para a classificação de ferimentos balísticos em imagens reais de cenas de crime da base de dados FDCPUnBGunshotDB. A arquitetura proposta utiliza uma rede convolucional pré-treinada (ResNet152) para extração de características espaciais robustas, substituindo o classificador denso tradicional por uma Kolmogorov-Arnold Network (KAN), capaz de modelar relações não-lineares complexas de forma mais eficiente.

A abordagem foi avaliada sob um esquema de validação cruzada k-fold rigoroso, prevenindo o vazamento de dados, e acompanhada de técnicas de Data Augmentation avançadas para mitigar o desbalanceamento inerente às amostras forenses. O desempenho do modelo TransferKAN foi comparado com resultados da literatura (artigo de 2024), demonstrando superioridade técnica.

Para a tarefa de classificação do tipo de ferimento (Entrada vs. Saída), a TransferKAN com validação k-fold obteve 86,0\% de Acurácia e uma AUC ponderada de 0,908, com F1-Score ponderado de 0,862. Na classificação da distância do disparo (MLSD), a TransferKAN apresentou uma melhoria em precisão ponderada, recall ponderado e F1-score balanceado, obtendo 87,3\% de Acurácia. Esses resultados evidenciam que a TransferKAN obteve a melhor performance e os melhores resultados gerais em métricas integradas, destacando o potencial da utilização das redes KAN como classificadores no auxílio de patologia forense em cenários reais.
